% ---------------------------------------------------------------------------
% NEW SECTION, drafted 2026-09-22. This is the material the abstract of
% tnf_paper.tex has always claimed (lines 122-131) and the body has never
% carried: the clause-folding census and the frequency-area law. Every number
% below is read out of the artefact tree and named with its file.
%
% TWO DISCREPANCIES AGAINST THE PUBLISHED ABSTRACT ARE MARKED \todo BELOW.
% They are stated rather than smoothed over; the author must settle which
% figure is correct before submission.
% ---------------------------------------------------------------------------

\section{Two measurements that were claimed and never shown}
\label{sec:themeasurements}

Both results in this section are about the instrument rather than about any
format. They are stated here because they are the conditions under which the
area and frequency figures of any comparison in this literature --- ours
included --- mean what they appear to mean, and because each was found by a
procedure a reader can run.

\subsection{Synthesis deletes the checks whose area is being measured}
\label{sec:folding}

On-die assertions are the usual defence against a design that routes and is
wrong. The defence has a failure mode that leaves every number intact: a clause
whose operands the optimiser can resolve at compile time is replaced by its
constant value, and a clause that has become a constant \emph{reads PASS in
every build, including the failing ones}.

The census in Table~\ref{tab:clauses} was taken with \texttt{tri clauses}, which
asks of each on-die clause whether it survives synthesis as logic. Across the
seven wrapper sources, \textbf{fourteen of twenty-eight clauses had folded to
constants} --- exactly half. One design (\texttt{gft\_signed\_dot4}) kept all
four; three designs kept only one each.

\begin{table}[htbp]\centering
\begin{tabular}{lcc}
\toprule
wrapper & clauses & folded \\
\midrule
\texttt{gft\_signed\_dot4\_jtag.v} & 4 & 0 \\
\texttt{gft\_signed\_mac\_jtag.v}  & 4 & 1 \\
\texttt{gft\_smul\_jtag.v}         & 4 & 2 \\
\texttt{gft\_dup\_jtag.v}          & 4 & 2 \\
\texttt{gft\_sadd\_jtag.v}         & 4 & 3 \\
\texttt{gft\_train1\_jtag.v}       & 4 & 3 \\
\texttt{gft\_xorpercep\_jtag.v}    & 4 & 3 \\
\midrule
total & 28 & 14 \\
\bottomrule
\end{tabular}
\caption{Clause census over the seven on-die wrappers, from
\texttt{measurements/clauses\_w983.json}. Half of the assertions that the
published verdicts rested on were constants at synthesis time.}
\label{tab:clauses}
\end{table}

% \todo THE ABSTRACT SAYS FIFTEEN. The shipped census says fourteen
% (totals: {wrappers: 7, clauses: 28, folded: 14, real: 14}). The census is
% named "census_of_todays_sources", so the abstract's figure may predate a
% repair. Settle this before submission; do not print both.

\paragraph{Three mechanisms, none of which announces itself.} Live counters were
already placed on the operands so that no clause could be satisfied by a folded
constant. The \emph{comparison} folds anyway, by three separate routes: a probe
register that nothing writes collapses to its \texttt{INIT}; two structurally
identical instances driven by identical operands are merged, so comparing them
is a tautology; and a clause both of whose operands are constants is evaluated
at compile time.

\paragraph{The obvious repair does not work.} Marking the probe register and the
instances \texttt{(* keep *)} is insufficient: \texttt{keep} preserves the cell,
and the optimiser still propagates the constant into the comparison. What worked
was making the sources \emph{structurally distinct while carrying equal values}
--- a second counter with the same seed and step feeding the control instance,
and a rotating probe register tested for a rotation-invariant property.

\paragraph{The size of the effect, measured directly.} On the repaired wrapper
\texttt{tri clauses} reports four clauses and none folded, and
\textbf{the LUT count rises from 452 to 696}. That difference is not an
overhead: it is the logic the folding had been deleting from every previous
build. \emph{Thirty-five per cent of the routed design was absent from the
measurements that were published as its area.}

\paragraph{Two published verdicts change meaning.} \texttt{gft\_sadd} was
published as ``PASS --- 4 of 4 clauses on the die''; of its four clauses three
were constants, so the honest reading is one of one measured clause passing.
\texttt{gft\_smul} was published as ``2 of 4 fail''; two of its four were
constants, so \emph{every clause that design actually evaluates on silicon is
false}. Neither verdict was wrong about the bits it reported. Both were wrong
about how many bits they were reporting.

\subsection{Frequency is a proxy for size, and the exponent survives the harness}
\label{sec:fmaxlaw}

If achieved frequency is largely determined by design size, then a metric of the
form ``throughput per unit area'' divides one quantity by a second that is
already a function of it, and a ranking by that metric restates a ranking by
size. Section~\ref{sec:selfcheck} reports the correlational form of this:
ranking by throughput per area reproduced ranking by $1/\text{area}$ at a
Spearman correlation of $0.961$. This subsection gives the functional form.

Eleven designs were placed and routed on an \texttt{xc7a200tfbg676-1} across
five placer seeds each, spanning $10$ to $396$ LUT --- a $39.6\times$ range of
size (Table~\ref{tab:fmaxrows}). Regressing the median achieved frequency on LUT
count in logarithms gives

\begin{equation}
F_{\max} \;=\; 4397\cdot \mathrm{LUT}^{-0.648},\qquad R^2 = 0.920,\quad n=11 .
\label{eq:fmaxlaw}
\end{equation}

Fitting instead on all fifty-five individual seed points, rather than on the
eleven medians, moves the exponent to $-0.646$ and $R^2$ to $0.921$: the law is
not an artefact of the median.

\begin{table}[htbp]\centering
\begin{tabular}{lrrr}
\toprule
design & LUT & $F_{\max}$ median (MHz) & seed spread \\
\midrule
\texttt{d\_baseline}  &  10 & 1004.0 &  7.9\% \\
\texttt{d\_int8}      &  10 &  995.0 &  5.6\% \\
\texttt{d\_bnf16}     &  12 &  693.0 & 16.9\% \\
\texttt{d\_gfternary} &  12 &  643.9 & 10.2\% \\
\texttt{d\_bin32}     &  19 &  985.2 &  5.7\% \\
\texttt{d\_tnf16}     &  19 &  719.4 & 17.9\% \\
\texttt{d\_bin16}     &  76 &  265.9 &  5.7\% \\
\texttt{d\_lns16}     & 132 &  304.9 & 20.5\% \\
\texttt{d\_posit16}   & 156 &  129.1 &  7.8\% \\
\texttt{d\_ibmhfp}    & 169 &  156.0 & 11.4\% \\
\texttt{d\_posit32}   & 396 &   74.2 &  6.2\% \\
\bottomrule
\end{tabular}
\caption{The eleven routed designs behind Equation~\ref{eq:fmaxlaw}, from
\texttt{measurements/decoder\_seed\_spread\_2026-08-19.json}. Five placer seeds
per design; the spread column is the seed-to-seed range of achieved frequency,
which reaches $20.5\%$ on one design.}
\label{tab:fmaxrows}
\end{table}

\paragraph{What makes this result stronger than it looks.} The harness behind
Table~\ref{tab:fmaxrows} observes only \texttt{q[7:0]\^{}q[31:24]}, and its own
record states that the \emph{absolute} costs are therefore cut, unequally, and
are not comparable between formats. A law fitted on cut absolute costs would
normally be worth nothing. But the cut multiplies LUT counts by a per-design
factor, which in logarithms is an additive per-design offset: it damages the
coefficient and leaves the exponent alone. That the exponent is recoverable from
a harness known to distort the absolute numbers is precisely the claim ---
$F_{\max}$ tracks size, and it does so whether or not the size was measured
faithfully.

The consequence for the literature is narrow and firm. A design that is smaller
will tend to clock faster for that reason alone; with a seed-to-seed spread
reaching $20.5\%$ on a single design, a ranking on frequency resolves only rows
separated by more than that; and a ratio of throughput to area is not an
independent figure of merit but a restatement of size. This does not invalidate
any frequency measurement. It invalidates the ratio built from them.

% \todo THE ABSTRACT PRINTS 3174 AND "a 41x area span". This fit reproduces
% the abstract's EXPONENT (-0.648) and R^2 (0.92) exactly, but gives 4397 and
% 39.6x. The published coefficient therefore came from a row set that is not
% pinned down by the artefacts in this tree. Either locate that set and ship
% it, or publish the fit above. Do not print 3174 without a file behind it.
